\documentclass[11pt,a4paper]{article}
\usepackage[T1]{fontenc}
\usepackage[utf8]{inputenc}
\usepackage[margin=27mm,headheight=15pt]{geometry}
\usepackage{amsmath,amssymb,amsthm,mathtools}
\usepackage{newpxtext,newpxmath}
\usepackage{microtype,booktabs,longtable,array}
\usepackage{xcolor,enumitem,fancyhdr}
\usepackage[most]{tcolorbox}
\usepackage{listings}
\usepackage{hyperref}
\usepackage[nameinlink,noabbrev]{cleveref}
\definecolor{ink}{HTML}{183D48}
\definecolor{accent}{HTML}{176B75}
\definecolor{pale}{HTML}{EDF5F5}
\definecolor{muted}{HTML}{59666B}
\hypersetup{colorlinks=true,linkcolor=accent,citecolor=accent,urlcolor=accent,
 pdftitle={Exactly One Digit per Height: Resolving Three Modular-Tetration Conjectures},
 pdfauthor={Computer-assisted research note}}
\pagestyle{fancy}\fancyhf{}
\fancyhead[L]{\small\color{muted}Exactly one digit per height}
\fancyhead[R]{\small\color{muted}Modular tetration / A324017}
\fancyfoot[C]{\small\thepage}
\renewcommand{\headrulewidth}{0.3pt}
\setlength{\parskip}{4pt}
\setlength{\parindent}{1.1em}
\setlength{\emergencystretch}{2em}
\numberwithin{equation}{section}
\theoremstyle{plain}
\newtheorem{theorem}{Theorem}[section]
\newtheorem{lemma}[theorem]{Lemma}
\newtheorem{corollary}[theorem]{Corollary}
\newtheorem{proposition}[theorem]{Proposition}
\theoremstyle{definition}
\newtheorem{definition}[theorem]{Definition}
\newtheorem{example}[theorem]{Example}
\newtheorem{remark}[theorem]{Remark}
\newcommand{\Z}{\mathbb Z}
\newcommand{\N}{\mathbb N}
\newcommand{\Q}{\mathbb Q}
\newcommand{\ordq}{\operatorname{ord}_q}
\newcommand{\vp}{v_p}
\newcommand{\TT}{T}
\newcommand{\eps}{\varepsilon}
\newcommand{\res}{\operatorname{res}}
\newtcolorbox{resultbox}[1][]{colback=pale,colframe=accent,boxrule=.6pt,
 arc=1.2mm,left=3mm,right=3mm,top=2mm,bottom=2mm,#1}
\lstset{language=Python,basicstyle=\ttfamily\small,keywordstyle=\color{accent}\bfseries,
 commentstyle=\color{muted},stringstyle=\color{ink},breaklines=true,
 frame=single,rulecolor=\color{accent},columns=fullflexible,keepspaces=true,
 showstringspaces=false,xleftmargin=1mm,xrightmargin=1mm}

\begin{document}
\thispagestyle{empty}
\begin{flushleft}
{\small\sffamily\color{accent}RESEARCH NOTE\quad /\quad 19 SEPTEMBER 2026}\par
\vspace{8mm}
{\LARGE\bfseries\color{ink}Exactly One Digit per Height}\par
\vspace{3mm}
{\Large Resolving three modular-tetration conjectures\par in OEIS A324017}\par
\vspace{5mm}
{\normalsize A computer-assisted research note}
\end{flushleft}
\vspace{3mm}
\begin{abstract}
OEIS A324017 records three conjectures about the residues of
$T_m=(2n-1)\uparrow\uparrow m$ modulo $(2n)^m$.  We prove the first two
and disprove the third for every nontrivial base in its stated domain.
The organizing result is stronger: if $q\ge4$ is even, $T_0=1$, and
$T_{h+1}=(q-1)^{T_h}$, then, for every $h\ge0$ and $r\ge1$,
\[
 T_{h+r}-T_h\equiv-2q^h\pmod{q^{h+1}}.
\]
Consequently, the tower acquires exactly one additional stable base-$q$
digit at each height.  We also determine every relevant prime-adic
valuation, prove uniqueness of the modular fixed point, give a
factorization-free digit-lifting algorithm, and extend the modular
stabilization result to arbitrarily high Knuth arrow ranks.  A local
Lambert $W$ representation explains the limiting values and the
essential sign distinction at the prime $2$.
All main proofs are elementary and self-contained.  A standard-library
Python implementation and independent exact-arithmetic checks accompany
the article.
\end{abstract}
\begin{resultbox}[title=Outcome and scope]
\textbf{Conjecture 1: true.}\quad
\textbf{Conjecture 2: true.}\quad
\textbf{Conjecture 3: false.}

The smallest allowed counterexample to the third statement is
\[
 3\uparrow\uparrow2=27\not\equiv59
   \equiv3\uparrow\uparrow3\pmod{64}.
\]
The contribution is an explicit resolution of the statements still
labelled conjectures in the retrieved OEIS entry, together with a
unified sharpening.  General modular stabilization and lifting-the-exponent
methods are established mathematics, not claimed discoveries here.
\end{resultbox}
\vfill
{\small\color{muted}Research status. The proofs have been checked algebraically and
against independent finite computations, but have not undergone independent
peer review. The literature check does not establish worldwide priority.}

\clearpage
\tableofcontents
\clearpage

\section{The selected problem and its status}\label{sec:question}

\subsection{The source and the three precise statements}
The chosen target is the set of conjectures attached to OEIS A324017,
an entry introduced by Davis Smith in 2019.  The entry retrieved on
19 September 2026 still labels all three statements as conjectures
\cite{oeis}.  This is a concrete, checkable source of the problem;
it is not a claim that no proof exists anywhere else.

Write
\begin{equation}\label{eq:setup}
 q=2n,\qquad a=q-1,\qquad T_0=1,\qquad T_{h+1}=a^{T_h}.
\end{equation}
Thus $T_h=a\uparrow\uparrow h$, with the exponentiation evaluated from
the top of the tower downward.  The array in the source is
\[
 A(m,n)=T_m\bmod q^m,
\]
where a residue always means the least nonnegative representative.
In this notation the three source conjectures are as follows.

\begin{description}[leftmargin=12mm,style=nextline]
\item[C1: fixed point.]
For $m,n,k>1$, if $T_m\equiv k\pmod{q^m}$, then
$a^k\equiv k\pmod{q^m}$.
\item[C2: second digit.]
For $n>1$ and $m\ge2$,
$\lfloor T_m/q\rfloor\equiv q-2\pmod q$.
\item[C3: early stabilization.]
For $m>1$, if $T_m\equiv j\pmod{q^{m+1}}$, then
$A(m+1,n)=j$, with $j$ understood as a canonical residue.
Equivalently, $T_{m+1}\equiv T_m\pmod{q^{m+1}}$.
\end{description}

The canonical-residue convention removes an irrelevant ambiguity in
C3: literal equality with an arbitrary congruent integer would already
be meaningless.  Even under its natural, most favorable interpretation,
C3 is false.  Our resolution concerns the statements as displayed;
we do not speculate about a different statement that the author might
have intended.

\subsection{Why this is an approachable research target}
The modulus is not arbitrary: the tower base is one less than an even
radix.  Every tower is odd, and
\[
 a\equiv-1\pmod q.
\]
These two observations neutralize the alternating sign in exponentiation.
Differences of tower heights can then be studied without constructing
any tower.  In particular, the difference of two exponentials factors
into a unit times an expression of the form $a^d-1$, with $d$ even.
That is exactly the setting in which elementary prime-adic valuations
are effective.

General stabilization of modular power towers and its algorithmic
implications are part of an existing literature.  Hittmeir studies the
relationship between general modular tetration and integer factorization
\cite{hittmeir}.  Rip\`a gives congruence-speed formulas for decimal
integer tetration \cite{ripa}.  Those works are important context:
proving that some digits eventually stabilize would not by itself be a
new contribution.  Here the target is the exact diagonal family
$q=a+1$, the exact height threshold, and the particular conjectures in
\cite{oeis}.

\subsection{A counterexample before the general theory}
Take $n=m=2$, so $q=4$ and $a=3$.  Then
\[
 T_2=27,\qquad T_3=3^{27}=7\,625\,597\,484\,987.
\]
Reducing modulo $q^{m+1}=64$ gives $27$ and $59$, respectively.
For a small independent certificate, note that
\[
 3^8\equiv33,\qquad 3^{16}\equiv1,\qquad
 3^{27}\equiv1\cdot33\cdot9\cdot3\equiv59\pmod{64}.
\]
Thus C3 fails at the first pair $m,n>1$ permitted by its nontrivial
interpretation.  This is not a rare anomaly: Theorem~\ref{thm:leading}
will show that it fails for every $n>1$ and every $m>1$.

\section{Notation and the meaning of an exact stable digit}

For a prime $p$ and a nonzero integer $x$, write $v_p(x)$ for the exponent
of $p$ in $x$.  When $q$ is composite, it is useful to introduce a
separate precision order
\begin{equation}\label{eq:ordq}
 \ordq(x)=\max\{s\ge0:q^s\mid x\}
 =\min_{p\mid q}\left\lfloor\frac{v_p(x)}{v_p(q)}\right\rfloor.
\end{equation}
This is not, in general, a multiplicative valuation.  For example,
$\operatorname{ord}_6(2)=\operatorname{ord}_6(3)=0$ but
$\operatorname{ord}_6(6)=1$.  We therefore use the phrase
\emph{base-$q$ precision order}, reserving ``valuation'' for $v_p$.

Two nonnegative integers have the same last $s$ base-$q$ digits exactly
when they are congruent modulo $q^s$.  If their precision order of
difference is $s$, then the digit in position $s$ differs, with the
units digit in position $0$.

The central claim can therefore be expressed without any numerical
approximation: every later tower agrees with $T_h$ in exactly its last
$h$ base-$q$ digits, and disagrees in the next digit.  The fact that
$q$ and the tower base are different is essential.  The phrase
``one digit per height'' always refers to radix $q$, not to radix $a$.

Throughout the main argument, $q\ge4$ is even, so $a=q-1\ge3$.
The excluded case $q=2$ has $a=1$ and $T_h=1$ for every height; it is
trivial but does not satisfy any nonzero-difference formula.

\section{Elementary valuation tools}\label{sec:lte}

We give the special lifting-the-exponent facts needed below, including
their proofs.  No theorem about real or complex tetration is involved.

\begin{lemma}\label{lem:principal}
Let $p$ be odd and let $X>1$ be an integer with $X\equiv1\pmod p$.  For every positive integer $d$,
\[
 v_p(X^d-1)=v_p(X-1)+v_p(d).
\]
The same formula holds for $p=2$ under the stronger assumption
$X\equiv1\pmod4$.
\end{lemma}
\begin{proof}
First suppose $p$ is odd and $s=v_p(X-1)\ge1$.
If $d$ is coprime to $p$, the factor
$1+X+\cdots+X^{d-1}$ is congruent to $d$ modulo $p$, so it is a unit.
Thus raising $X$ to a $p$-coprime exponent leaves the valuation of
$X-1$ unchanged.

Writing $X=1+p^su$, with $p\nmid u$, the binomial expansion of $X^p-1$
has first term $p^{s+1}u$.  Every other term has valuation at least
$s+2$.  For the last term this uses $ps\ge s+2$, which holds for
$p\ge3$ and $s\ge1$.  Consequently, raising to the $p$th power adds
exactly one to the valuation.  Factor $d=p^tw$, with $p\nmid w$, and
apply these two observations successively.

For $p=2$, odd exponents are handled by the same geometric sum.
If $X\equiv1\pmod4$, then $v_2(X+1)=1$, whence
$v_2(X^2-1)=v_2(X-1)+1$.  Repeated squaring and the odd part of $d$
complete the proof.
\end{proof}

Set
\begin{equation}\label{eq:uvc}
 u=v_2(q-2)=v_2(a-1),\qquad v=v_2(q)=v_2(a+1),\qquad c=u+v-1.
\end{equation}
Both $u$ and $v$ are positive, one of them is $1$, and $c\ge v$.

\begin{lemma}\label{lem:even-exponent}
For every positive even integer $d$ and every odd prime $p\mid q$,
\begin{align}
 v_p(a^d-1)&=v_p(q)+v_p(d),\label{eq:odd-lte}\\
 v_2(a^d-1)&=c+v_2(d).\label{eq:two-lte}
\end{align}
If $t$ is a positive odd integer, then $v_2(a^t-1)=u$.
\end{lemma}
\begin{proof}
Write $a^d=(a^2)^{d/2}$.  For odd $p\mid q$, the identity
$a^2-1=(q-2)q$ gives $v_p(a^2-1)=v_p(q)$, because
$p\nmid q-2$.  Apply Lemma~\ref{lem:principal}; division of $d$ by $2$
does not affect its $p$-adic valuation.

At $p=2$, $a^2\equiv1\pmod8$ and $v_2(a^2-1)=u+v$.
The same lemma gives $u+v+v_2(d/2)=c+v_2(d)$.
For odd $t$, the geometric sum in
$a^t-1=(a-1)(1+a+\cdots+a^{t-1})$ is odd.
\end{proof}

\section{Exact distance between any two tower heights}\label{sec:distance}

\begin{theorem}[Exact local distances]\label{thm:local}
For every even $q\ge4$, every $h\ge0$, every $r\ge1$, and every odd
prime $p\mid q$,
\begin{align}
 v_p(T_{h+r}-T_h)&=h\,v_p(q),\label{eq:mainodd}\\
 v_2(T_{h+r}-T_h)&=u+h(u+v-1).\label{eq:maintwo}
\end{align}
In particular, the answer is independent of the positive height gap $r$.
\end{theorem}
\begin{proof}
Fix $r$ and put $D_h=T_{h+r}-T_h$.  Since $a\ge3$, the towers strictly
increase, so all $D_h$ are positive.  Every tower is odd, making every
$D_h$ even.

At height $0$, $D_0=T_r-1=a^{T_{r-1}}-1$.
For an odd prime $p\mid q$, the oddness of $T_{r-1}$ gives
$T_r\equiv-1\pmod p$, and hence $D_0\equiv-2\pmod p$.
Thus $v_p(D_0)=0$.  The last assertion of
Lemma~\ref{lem:even-exponent} gives $v_2(D_0)=u$.

For the induction step, factor
\begin{equation}\label{eq:Drec}
 D_{h+1}=a^{T_h}\bigl(a^{D_h}-1\bigr).
\end{equation}
The first factor is a unit at every prime dividing $q$.
Since $D_h$ is even, Lemma~\ref{lem:even-exponent} adds $v_p(q)$ at each
odd prime and adds $c$ at $2$.  Starting from the height-zero values
gives exactly \eqref{eq:mainodd} and \eqref{eq:maintwo}.
\end{proof}

For clarity, the $2$-adic formula has the following two forms:
\begin{equation}\label{eq:two-cases}
 v_2(T_{h+r}-T_h)=
 \begin{cases}
 1+h\,v_2(q),&q\equiv0\pmod4,\\
 (h+1)v_2(q-2),&q\equiv2\pmod4.
 \end{cases}
\end{equation}

\begin{corollary}[Exactly one radix digit per height]\label{cor:order}
For every $h\ge0$ and $r\ge1$,
\begin{equation}\label{eq:exact-q}
 \ordq(T_{h+r}-T_h)=h.
\end{equation}
Moreover, $2q^h\mid T_{h+r}-T_h$.
\end{corollary}
\begin{proof}
The local formulas imply $q^h\mid D_h$.  At $2$ they give the stronger
inequality $v_2(D_h)=u+hc\ge1+hv$, proving $2q^h\mid D_h$.

If $q$ has an odd prime factor $p$, then
$v_p(D_h)=h\,v_p(q)<(h+1)v_p(q)$, so $q^{h+1}\nmid D_h$.
If $q$ has no odd prime factor, write $q=2^v$ with $v\ge2$.
Then $u=1$ and $v_2(D_h)=1+hv<(h+1)v$.
This proves \eqref{eq:exact-q} in all cases.
\end{proof}

\begin{corollary}[Sharp stabilization threshold]\label{cor:threshold}
Fix $m\ge1$.  The residues $T_h\bmod q^m$ become constant precisely
at height $m$.  More strongly, no height $h<m$ has the eventual residue.
\end{corollary}
\begin{proof}
If $h\ge m$, \eqref{eq:exact-q} shows that every $T_{h+r}$ agrees
with $T_h$ modulo $q^m$.  If $h<m$, then $T_m-T_h$ has precision
order $h<m$, so $T_h$ differs from the eventual residue $T_m\bmod q^m$.
\end{proof}

This is stronger than saying that neighboring residues eventually agree:
it rules out an earlier accidental visit to the final residue followed
by a departure.  The exact formulas also explain why the excess powers
of $2$ do not speed up the full radix-$q$ stabilization.  An odd prime
factor is a bottleneck unless $q$ is a power of $2$; in that exceptional
case the single extra factor of $2$ is still less than an entire new
radix digit.

\section{The universal first unstable-digit difference}\label{sec:leading}

The preceding result determines the first location of disagreement.
We now determine the difference at that location.  A factor of $2$ in
Corollary~\ref{cor:order} is crucial when the radix is even.

\begin{lemma}[A binomial precision estimate]\label{lem:binomial}
Let $q\ge4$ be even, let $h\ge0$, and let $d>0$ be divisible by
$2q^h$.  Then
\begin{equation}\label{eq:linearized}
 (q-1)^d\equiv1-qd\pmod{q^{h+2}}.
\end{equation}
\end{lemma}
\begin{proof}
The exponent $d$ is even, so $(q-1)^d=(1-q)^d$.
We must show that every binomial term
$\binom dj(-q)^j$ with $j\ge2$ is divisible by $q^{h+2}$.
For a prime $p\mid q$, put $e=v_p(q)$.  The identity
\[
 \binom dj=\frac dj\binom{d-1}{j-1}
\]
shows that the valuation of this term is at least
\[
 je+v_p(d)-v_p(j).
\]
For odd $p$, this is at least $(h+j)e-v_p(j)$, which is at least
$(h+2)e$ because $v_p(j)\le j-2$ for $j\ge2$.
For $p=2$, the extra factor $2$ in $d$ gives the lower bound
$(h+j)e+1-v_2(j)$.  It is at least $(h+2)e$, since
$v_2(j)\le j-1$ and $e\ge1$.
This holds at every prime dividing $q$, proving the claimed divisibility.
\end{proof}

\begin{theorem}[Universal leading difference]\label{thm:leading}
For every even $q\ge4$, every $h\ge0$, and every $r\ge1$,
\begin{equation}\label{eq:leading}
 \boxed{\ T_{h+r}-T_h\equiv-2q^h\pmod{q^{h+1}}\ }.
\end{equation}
Equivalently,
\begin{equation}\label{eq:normalized}
 \frac{T_{h+r}-T_h}{q^h}\equiv q-2\pmod q.
\end{equation}
\end{theorem}
\begin{proof}
Again fix $r$ and put $D_h=T_{h+r}-T_h$.
At height zero, $D_0\equiv-2\pmod q$, because every positive-height
tower is congruent to $-1$ modulo $q$.

By Corollary~\ref{cor:order}, $2q^h\mid D_h$.
Apply Lemma~\ref{lem:binomial} to \eqref{eq:Drec}:
\[
 D_{h+1}\equiv a^{T_h}(-qD_h)\pmod{q^{h+2}}.
\]
Since $T_h$ is odd, $a^{T_h}\equiv-1\pmod q$.
Also, $qD_h$ is divisible by $q^{h+1}$.  Replacing the first factor
by $-1$ is therefore legitimate modulo $q^{h+2}$, and gives
\[
 D_{h+1}\equiv qD_h\pmod{q^{h+2}}.
\]
Induction from $D_0\equiv-2\pmod q$ proves \eqref{eq:leading}.
\end{proof}

\begin{corollary}[Complete refutation and correction of C3]\label{cor:C3}
For every $n>1$ and $m>1$, C3 is false.  Its correct replacement is
\begin{equation}\label{eq:correct-C3}
 A(m+1,n)\equiv T_m-2(2n)^m\pmod{(2n)^{m+1}}.
\end{equation}
For $n=1$, the tower is constantly $1$ and the original C3 is trivially true.
\end{corollary}
\begin{proof}
Use \eqref{eq:leading} with $h=m$ and $r=1$.
The residue $-2q^m$ is nonzero modulo $q^{m+1}$ because $q\ge4$ does
not divide $2$.  The corrected congruence is the same identity rearranged.
\end{proof}

\begin{remark}
The number $q-2$ is the \emph{difference} of the first unstable digits,
modulo $q$.  It is not a claim that the digit itself always equals
$q-2$.  Digits farther to the left can also change, and carries must
not be ignored.
\end{remark}

\section{Modular fixed points and the two true conjectures}\label{sec:fixed}

\subsection{Why exponent reduction is valid here}
Blindly replacing a tower exponent by its residue modulo the output
modulus is invalid in general.  For example,
$2^3\not\equiv2^{3+3}\pmod3$.
Our restricted family admits precisely the compatibility that such a
computation needs.

\begin{lemma}[Exponent period and contraction]\label{lem:contraction}
For every $s\ge1$,
\begin{equation}\label{eq:period}
 a^{q^s}\equiv1\pmod{q^{s+1}}.
\end{equation}
If $x,y$ are nonnegative odd integers and $q^s\mid x-y$, where now
$s\ge0$, then
\begin{equation}\label{eq:contraction}
 a^x\equiv a^y\pmod{q^{s+1}}.
\end{equation}
\end{lemma}
\begin{proof}
For \eqref{eq:period}, the exponent $q^s$ is even.
The valuation formulas in Lemma~\ref{lem:even-exponent} give
$(s+1)v_p(q)$ at each odd prime and at least $(s+1)v_2(q)$ at $2$.

For the second statement, interchange $x,y$ if necessary and factor
$a^x-a^y=a^y(a^{x-y}-1)$.  If $x=y$ there is nothing to prove.
Otherwise $x-y$ is even.  At an odd prime, exponentiation adds
$v_p(q)$ to the valuation of the difference; at $2$ it adds
$c\ge v_2(q)$.  The case $s=0$ is included in this argument.
\end{proof}

The period identity shows that
\[
 F_m:\Z/q^m\Z\longrightarrow\Z/q^m\Z,\qquad x\longmapsto a^x
\]
is well defined, when evaluated on nonnegative representatives.
It always returns an odd residue.  The contraction identity says that,
on odd residues, one application gains at least one radix digit of
agreement.

\begin{theorem}[Unique modular fixed point]\label{thm:fixed}
For each $m\ge1$, the congruence
\begin{equation}\label{eq:fixed}
 a^x\equiv x\pmod{q^m}
\end{equation}
has exactly one residue-class solution.  It is
\[
 R_m:=T_m\bmod q^m.
\]
The representatives satisfy
\begin{equation}\label{eq:compatibility}
 R_{m+1}\equiv R_m\pmod{q^m}.
\end{equation}
\end{theorem}
\begin{proof}
By the period property, $a^{R_m}\equiv a^{T_m}=T_{m+1}\pmod{q^m}$.
By Corollary~\ref{cor:threshold}, $T_{m+1}\equiv T_m\equiv R_m\pmod{q^m}$.
This proves existence.

If $x$ and $y$ are two fixed points, both are odd.
Starting with their agreement modulo $q^0$, apply
\eqref{eq:contraction} and the fixed-point equations to obtain their
agreement modulo $q$, then modulo $q^2$, and so on through $q^m$.
This proves uniqueness.  Compatibility follows directly from the
stabilization of the towers.
\end{proof}

\begin{corollary}[C1 is true]\label{cor:C1}
For every $m\ge1$, every $n>1$, and every nonnegative integer $k$
with $T_m\equiv k\pmod{(2n)^m}$,
\[
 (2n-1)^k\equiv k\pmod{(2n)^m}.
\]
Thus the source's restrictions $m,k>1$ are unnecessary.
\end{corollary}
\begin{proof}
Such a $k$ represents the unique fixed-point class $R_m$.
The exponent period ensures that changing its representative does not
change the value of $a^k$ modulo $q^m$.
\end{proof}

\begin{theorem}[Explicit two-digit residue; C2 is true]\label{thm:C2}
For $n>1$ and every $h\ge2$,
\begin{equation}\label{eq:twodigit}
 T_h\equiv q^2-q-1\pmod{q^2}.
\end{equation}
In particular, $\lfloor T_h/q\rfloor\equiv q-2\pmod q$.
\end{theorem}
\begin{proof}
The exponent $a=q-1$ is odd.  A first-order binomial expansion gives
\[
 T_2=(q-1)^{q-1}=-(1-q)^{q-1}
 \equiv-1+q(q-1)\equiv-1-q\pmod{q^2}.
\]
Its canonical representative is $q^2-q-1$.
All heights $h\ge2$ have the same residue modulo $q^2$ by
Corollary~\ref{cor:threshold}.
Finally,
\[
 q^2-q-1=q(q-2)+(q-1).
\]
Dividing $T_h$ by $q$ and reducing the integer quotient modulo $q$
therefore gives $q-2$.
\end{proof}

\section{The limiting integer and its digits}\label{sec:limit}

\subsection{An inverse limit, not a real infinite tower}
The compatible residues $R_m$ define an element
\[
 \xi_q\in\Z_q:=\varprojlim_m\Z/q^m\Z.
\]
For composite $q$, this is a complete ring rather than a field.
By the Chinese remainder theorem, it can be identified with
$\prod_{p\mid q}\Z_p$.  We do not interpret $\xi_q$ as a finite real
limit: since $a\ge3$, the ordinary real tower diverges to infinity.
Only its modular tails converge.

\begin{proposition}\label{prop:limit-error}
For every $h\ge0$,
\begin{equation}\label{eq:limit-leading}
 \xi_q-T_h\equiv-2q^h\pmod{q^{h+1}},
 \qquad \ordq(\xi_q-T_h)=h.
\end{equation}
Each local component $\xi_{q,p}$ also satisfies the exact valuations
in Theorem~\ref{thm:local}, with $T_{h+r}$ replaced by $\xi_{q,p}$.
\end{proposition}
\begin{proof}
Every sufficiently late tower represents $\xi_q$ at any fixed
precision.  Thus the congruence in Theorem~\ref{thm:leading} passes directly
to the inverse limit.  For a prime $p\mid q$, use precision one higher
than the valuation in Theorem~\ref{thm:local}.  All later differences are
nonzero modulo that prime power, and convergence makes their residue
eventually constant; its valuation is therefore unchanged.
\end{proof}

\subsection{Digit lifting without building towers}
The next theorem gives a simple recurrence for the array itself.

\begin{theorem}[Canonical digit lift]\label{thm:lift}
The fixed-point representatives are obtained by
\begin{equation}\label{eq:lift}
 R_1=q-1,\qquad R_{m+1}=a^{R_m}\bmod q^{m+1}.
\end{equation}
Equivalently, if $R_{m+1}=R_m+d_mq^m$ with $0\le d_m<q$, then
\begin{equation}\label{eq:newdigit}
 d_m\equiv\frac{a^{R_m}-R_m}{q^m}\pmod q.
\end{equation}
The numerator in the quotient is divisible by $q^m$.
\end{theorem}
\begin{proof}
Since $T_m\equiv R_m\pmod{q^m}$ and both are odd,
Lemma~\ref{lem:contraction} gives
$a^{T_m}\equiv a^{R_m}\pmod{q^{m+1}}$.
The left side is $T_{m+1}$, proving \eqref{eq:lift}.
Compatibility gives \eqref{eq:newdigit}; the fixed-point property at
precision $m$ proves the divisibility of its numerator.
\end{proof}

\begin{table}[ht]
\centering
\caption{Computed stable representatives $R_m$. The modulus in each column is
$q^m$, so all entries are ordinary decimal displays of different-radix residues.}
\label{tab:residues}
\begin{tabular}{r r r r}
\toprule
$m$ & $q=4$, base $3$ & $q=6$, base $5$ & $q=10$, base $9$\\
\midrule
1&3&5&9\\
2&11&29&89\\
3&59&29&289\\
4&59&1109&5289\\
5&827&3701&45289\\
6&2875&34805&745289\\
7&15163&128117&2745289\\
8&31547&687989&92745289\\
\bottomrule
\end{tabular}
\end{table}

The repeated value $R_3=R_4=59$ at $q=4$ is not a counterexample to
exact one-digit stabilization.  It says that one digit of the limit is
zero.  It does \emph{not} say that $T_3$ is already correct modulo $4^4$:
in fact $T_3\equiv187\pmod{256}$, whereas the stable residue is $59$.
Confusing an already truncated representative $R_m$ with the full tower
$T_m$ would conceal exactly the discrepancy that refutes C3.

\section{Algorithms and independent exact verification}\label{sec:algorithms}

\subsection{A factorization-free evaluator}
The following is the essential stable-residue algorithm, reproduced in
the accompanying \texttt{code/tower\_digits.py} with input validation.
\begin{lstlisting}
def stable_residue(q, precision):
    a, modulus, x = q - 1, q, q - 1
    for _ in range(2, precision + 1):
        modulus *= q
        x = pow(a, x, modulus)
    return x
\end{lstlisting}

To evaluate $T_h\bmod q^m$ for a finite height, first replace $h$ by
$\min(h,m)$.  At a fixed modulus $M=q^m$, iterate
$x\leftarrow a^x\bmod M$ from $x=1$ that many times.
The period in Lemma~\ref{lem:contraction} justifies reducing each exponent
modulo $M$; Corollary~\ref{cor:threshold} justifies the height cap.
The implementation uses digit lifting for the stabilized case and the
fixed-modulus loop for smaller heights.

Let $L=\lfloor\log_2(q^m)\rfloor+1$ be the output-modulus bit length.
There are at most $m$ modular exponentiations, and every exponent stored
has at most $L$ bits.  With ordinary schoolbook multiplication and
reduction, square-and-multiply gives the safe bit-complexity bound
\[
 O(mL^3)
\]
and $O(L)$ working storage for a single evaluation, apart from optional
caching of completed queries.  This is a deliberately conservative
bound, not a claim of optimality.  Crucially, the computation does not
factor $q$, construct $T_h$, or depend on a tower-sized exponent.
The restriction $a=q-1$ is why this does not provide an unrestricted
fast modular-tetration or factoring algorithm.

\subsection{An independent oracle}
A test based solely on the same fixed-modulus iteration could reproduce
the same mistake as the implementation.  The verification script therefore
uses a second method: a totient chain with exact capped exponents.
For a general modulus $M$, put $P=\varphi(M)$.  If the true exponent
$E$ is less than $P$, it is computed exactly.  Otherwise it is replaced by
\[
 (E\bmod P)+P.
\]
Why is this safe even for a nonunit base?  At a prime power dividing $M$
where the base is a unit, the relevant multiplicative period divides
$P$.  At a prime power $p^e$ where the base is divisible by $p$, both
exponents are at least $P\ge\varphi(p^e)\ge e$, so both powers vanish
modulo $p^e$.  The Chinese remainder theorem finishes the justification.

The oracle is itself checked against directly constructed, manageable
integer towers at thousands of small parameter values.  Thus the
validation has three levels: literal integers, a general modular oracle,
and the specialized algorithm.

\subsection{Recorded checks}
The supplied run completed \textbf{62,649 exact assertions}.  It covered
every even $q$ from $4$ through $200$, heights $0$ through $10$, gaps
$1,2,3$, and output precisions through $10$, in the combinations detailed
in \texttt{data/verification.json}.  It also examined all
\textbf{51,414 residue candidates} in a separate exhaustive fixed-point
search for even $q\le24$ and precisions $1,2,3$.

\begin{table}[ht]
\centering
\caption{Selected assertion counts from the reproducible test report.}
\begin{tabular}{l r}
\toprule
Check family & Assertions\\
\midrule
General oracle against literal integers&8700\\
Universal leading-difference congruence&3267\\
Exact prime-adic distances&7260\\
Specialized tower residues against the oracle&10890\\
Fixed-point representatives&2970\\
C3 fails throughout the tested nontrivial range&891\\
Capped hyperoperations against small exact values&21000\\
Local Lambert-series residues&60\\
\bottomrule
\end{tabular}
\end{table}

All these checks use integer or rational arithmetic; no floating-point
agreement is being mistaken for proof.  The finite computations are
supporting evidence and software validation.  The infinite statements
follow from the preceding proofs, not from extrapolation from a search.

\section{From tetration to higher Knuth arrows}\label{sec:arrows}

The selected problem lies directly in tetration, but the same result
also makes certain higher hyperoperations tractable modulo $q^m$.
We state all normalization conventions to avoid an arrow-count shift.

For $a=q-1\ge3$, define
\begin{align*}
 H_1(n)&=a^n,\\
 H_r(0)&=1 &&(r\ge2),\\
 H_r(n+1)&=H_{r-1}(H_r(n)) &&(r\ge2).
\end{align*}
Thus $H_r(n)=a\uparrow^r n$ and $H_2(n)=T_n$.

\begin{lemma}[Every higher operation is an integer-height tower]
\label{lem:height}
There are integers $J_r(n)\ge0$, for $r\ge2$, such that
\[
 H_r(n)=T_{J_r(n)}.
\]
They can be specified by
\begin{align}
 J_2(n)&=n,\label{eq:Jbase}\\
 J_r(0)&=0,\qquad
 J_r(n+1)=J_{r-1}(H_r(n))\quad(r\ge3).\label{eq:Jrec}
\end{align}
Furthermore, $J_r(1)=1$ and
\begin{equation}\label{eq:Jbound}
 J_r(n)\ge n+r-2\qquad(r\ge2,\ n\ge2).
\end{equation}
\end{lemma}
\begin{proof}
For $r=2$ the representation is the definition of tetration.
Assuming it at rank $r-1$,
\[
 H_r(n+1)=H_{r-1}(H_r(n))
         =T_{J_{r-1}(H_r(n))},
\]
which proves the representation at rank $r$.
The identities at $n=0,1$ follow immediately.

An induction on rank and argument gives
$H_r(n)\ge a^n$ for every $r\ge1$ and $n\ge0$.
Indeed, at rank $r$ the recursion yields
$H_r(n+1)\ge a^{H_r(n)}\ge a^{n+1}$.
For $n\ge2$ and $a\ge3$, one has $a^{n-1}\ge n+1$.
Now induct on $r$ in \eqref{eq:Jbound}.  The case $r=2$ is equality.
At rank $r\ge3$, the argument $H_r(n-1)$ is at least $n+1\ge3$, so
\[
 J_r(n)=J_{r-1}(H_r(n-1))
 \ge H_r(n-1)+r-3\ge n+r-2.
\]
\end{proof}

\begin{theorem}[Modular saturation across arrow ranks]\label{thm:arrows}
For $r\ge2$ and $n\ge0$,
\begin{equation}\label{eq:arrowiff}
 H_r(n)\equiv R_m\pmod{q^m}
 \quad\Longleftrightarrow\quad J_r(n)\ge m.
\end{equation}
In particular, for $n\ge2$ the readily checkable condition
\begin{equation}\label{eq:arrowsufficient}
 n+r-2\ge m
\end{equation}
is sufficient.  Thus, at fixed precision, sufficiently high arrow
ranks give exactly the same modular tail as a sufficiently tall tower.
\end{theorem}
\begin{proof}
The equivalence follows from the representation in Lemma~\ref{lem:height}
and the sharp threshold in Corollary~\ref{cor:threshold}.  The sufficient
condition follows from \eqref{eq:Jbound}.
\end{proof}

For example, the supplied implementation evaluates
$9\uparrow^{1\,000\,000}2$ modulo $10^{60}$ as
\begin{equation*}
 \begin{split}
 &\texttt{864868894914047889985007525482}\\[-1mm]
 &\hspace{9mm}\texttt{726503475610748087597392745289},
 \end{split}
\end{equation*}
where the two lines concatenate to the $60$-digit decimal
residue.  This result is simply
$R_{60}$ for $q=10$; the million-arrow expression is never expanded.

For queries below the sufficient saturation bound, the code uses exact
capped hyperoperations.  To compute $\min(H_r(n),C)$, let $t$ be the
least integer with $a^t\ge C$.  The inequalities
\[
 H_r(n)\ge a^n,\qquad H_r(n)\ge a^r\quad(r\ge2,\ n\ge2)
\]
allow immediate return of $C$ whenever $n\ge t$ or $r\ge t$,
after handling $n=0,1$.  Otherwise the small recursion can be evaluated
without creating any integer exceeding the cap, apart from bounded
intermediate products used to determine the threshold.
In a modular query the needed cap is only $m$.

The assumption $a\ge3$ matters.  One must not extend the rank bound to
$a=2$: $2\uparrow^r2=4$ for every positive rank $r$.
Nor does Theorem~\ref{thm:arrows} assert that arbitrary Conway chains or
arbitrary fast-growing-hierarchy notations have been evaluated.  It is
a precise extension to Knuth hyperoperations with the stated base and
modulus relationship.

\section{A local Lambert \texorpdfstring{$W$}{W} description}\label{sec:lambert}

This optional section expresses the limiting residue system using
$p$-adic analysis.  It is not needed for any of the three conjectures.
The $p$-adic Lambert function and its power series are established
objects; see Mez\H{o} \cite{mezo}.  Our task is to select the correct
local sign for this tower family.

\subsection{Principal units and the sign at two}
Fix a prime $p\mid q$ and define
\begin{equation}\label{eq:epsilon}
 \eps_p=
 \begin{cases}
 +1,&p=2\text{ and }q\equiv2\pmod4,\\
 -1,&\text{otherwise},
 \end{cases}
 \qquad b_p=\eps_p(q-1),\qquad \ell_p=\log_p(b_p).
\end{equation}
For odd $p$, $b_p=1-q\in1+p\Z_p$.
For $p=2$, the definition ensures $b_p\in1+4\Z_2$.
Consequently logarithm and exponential are inverse on the relevant
principal-unit neighborhoods.  More explicitly,
\[
 v_p(\ell_p)=
 \begin{cases}
 v_p(q),&p\text{ odd},\\
 v_2(q),&p=2,\ q\equiv0\pmod4,\\
 v_2(q-2),&p=2,\ q\equiv2\pmod4.
 \end{cases}
\]
The logarithm is nonzero: $b_p$ is neither the identity nor a
nontrivial torsion element in these principal-unit neighborhoods.

For an ordinary odd integer $x$,
\begin{equation}\label{eq:localmap}
 a^x=\eps_p\exp_p(\ell_p x).
\end{equation}
The tower exponents are all odd, so this is exactly the analytic map
that must be continued to the limiting component.  At odd $p$ the sign
$-1$ is prescribed by the odd tower branch; it should not be reconstructed
from a nonexistent continuous parity function on $\Z_p$.

\subsection{The closed formula and its proof}
Define the small-argument Lambert series
\begin{equation}\label{eq:Wseries}
 W_p(z)=\sum_{j\ge1}\frac{(-j)^{j-1}}{j!}z^j.
\end{equation}
Formal Lagrange inversion of $w\mapsto we^w$ gives these coefficients:
$[z^j]W(z)=j^{-1}[t^{j-1}]e^{-jt}=(-j)^{j-1}/j!$.
The series converges when $v_p(z)>1/(p-1)$.  For the present arguments,
this follows directly from
\[
 v_p\!\left(\frac{(-j)^{j-1}}{j!}z^j\right)
 \ge jv_p(z)-\frac{j-1}{p-1}\longrightarrow+\infty.
\]
The inverse identity $W_p(z)\exp_p(W_p(z))=z$ therefore holds in
this neighborhood.  These facts agree with the standard $p$-adic
Lambert theory in \cite{mezo}.

\begin{theorem}[Local limit formula]\label{thm:W}
For every prime $p\mid q$, the component of the stable tower limit is
\begin{equation}\label{eq:Wanswer}
 \boxed{\ \xi_{q,p}=-\frac{W_p(-\eps_p\ell_p)}{\ell_p}\ }.
\end{equation}
All series in this expression converge in $\Q_p$.
\end{theorem}
\begin{proof}
The valuations of $\ell_p$ are at least $1$ for odd $p$ and at least
$2$ for $p=2$, so $v_p(\ell_p)>1/(p-1)$.
By the exact local distances, $T_h$ converges to $\xi_{q,p}$.
Continuity in \eqref{eq:localmap} gives
\[
 \xi_{q,p}=\eps_p\exp_p(\ell_p\xi_{q,p}).
\]
Put $w=-\ell_p\xi_{q,p}$.  Then
$w\exp_p(w)=-\eps_p\ell_p$.
The small inverse branch is unique: on its domain the linear term of
$w\exp_p(w)$ preserves distances and the higher terms are strictly
smaller.  Hence $w=W_p(-\eps_p\ell_p)$, proving the formula.
Alternatively, substituting the right side of \eqref{eq:Wanswer} into
the fixed-point equation and using the strict contraction of
$x\mapsto\eps_p\exp_p(\ell_p x)$ proves the same identification.
\end{proof}

\begin{example}[Why a uniform sign would be wrong]
Take $q=6$, so $a=5$.  The $2$-adic component uses
$\eps_2=1$ and $\ell_2=\log_2(5)$; it is congruent to $1$ modulo $4$.
The $3$-adic component uses $\eps_3=-1$ and
$\ell_3=\log_3(-5)$; it is congruent to $-1$ modulo $3$.
A blind use of the minus-sign formula at both primes would give the
wrong $2$-adic branch, congruent to $-1$ modulo $4$.
Extending the $2$-adic logarithm by $\log_2(-1)=0$ does not repair the
mistake: logarithm then forgets exactly the sign that exponentiation
of an odd integer retains.
\end{example}

\subsection{Reproducible rational-series verification}
The function \texttt{local\_lambert\_residue} evaluates finite logarithm
and Lambert polynomials using exact rational numbers.  It is intended
as an independent educational check, not as the fastest algorithm.
At a requested $p$-adic precision $K$, it sets an auxiliary target of
$K+4$ digits and uses $2(K+4)+4$ terms in each series.

Here are sufficient tail estimates.  If $s=v_p(b_p-1)$, a logarithm
term of degree $j$ has valuation $js-v_p(j)$.  Since
$v_p(j)\le j/2$ for $j\ge2$, the chosen truncation makes the entire
logarithm tail vanish beyond the requested precision.
For the fixed-point series obtained by dividing \eqref{eq:Wseries}
by $\ell_p$, a term with index $j$ has valuation at least
\[
 (j-1)s-v_p(j!).
\]
For odd $p$, this is at least $(j-1)/2$; for $p=2$, where $s\ge2$,
it is at least $j-1$.  These bounds justify the chosen truncation.
Substituting the truncated logarithm preserves the target precision:
the linear coefficient of the fixed-point series is $1$, and the
higher powers on these principal-unit domains cannot magnify an input
error.  This can also be checked termwise using
$x^k-y^k=(x-y)\sum_{i=0}^{k-1}x^{k-1-i}y^i$ and the same factorial
bounds.

The supplied checks compare this rational-series evaluation with the
general modular oracle at $60$ local parameter triples, including
both residue classes of $q$ modulo $4$.

\section{Conclusions and the remaining boundary of the result}

The source problem has a complete resolution in its stated domain.
C1 is true, and in fact there is a unique fixed-point residue modulo
every $q^m$.  C2 is true, with the stronger explicit formula
$T_h\equiv q^2-q-1\pmod{q^2}$ for all $h\ge2$.
C3 is false for every $q=2n\ge4$; the error is not merely nonzero but
always $-2q^m$ modulo $q^{m+1}$.

The unified theorem is
\[
 T_{h+r}-T_h\equiv-2q^h\pmod{q^{h+1}},
 \qquad h\ge0,\quad r\ge1.
\]
It yields an exact stabilization threshold, a digit-lifting recurrence,
and modular saturation across sufficiently large Knuth arrow ranks.
The prime-adic analysis explains why the factor $2$ in the leading
error is universal and why it never supplies a full additional radix
digit when $q\ge4$.

Several possible continuations are natural, but are not being labelled
verified open problems here.  One could classify similarly explicit
first-unstable-digit laws when the base is congruent to a root of unity
other than $-1$; determine sharper sufficient rank bounds than
$n+r-2\ge m$; or build a finite-precision evaluator for broader classes
of chained-arrow expressions.  Each would need its own precise source
check and proof.  Nothing in this note resolves the general analytic
extension of tetration, arbitrary-modulus hyperoperation complexity,
or the comparison of unrestricted ordinal-indexed fast-growing
hierarchies.

\appendix
\section{A compact replacement for the OEIS conjectural comments}
The following mathematical text is suitable as a proposed correction;
it has not been submitted to OEIS by this work.

\begin{resultbox}
Let $q=2n\ge4$, $T_0=1$, and $T_{h+1}=(q-1)^{T_h}$.
For all $h\ge0$ and $r\ge1$,
\[
 T_{h+r}-T_h\equiv-2q^h\pmod{q^{h+1}}.
\]
Thus the first two conjectures are true.  In fact, $A(m,n)$ is the
unique solution of $(q-1)^x\equiv x\pmod{q^m}$, and
$T_h\equiv q^2-q-1\pmod{q^2}$ for $h\ge2$.
The third conjecture is false for every $n>1$: the correct relation is
$A(m+1,n)\equiv T_m-2q^m\pmod{q^{m+1}}$.
For example, $3\uparrow\uparrow2\equiv27$ but
$3\uparrow\uparrow3\equiv59\pmod{64}$.
\end{resultbox}

\section{Proof-audit checklist}
A reader can audit the principal claims without examining the entire
software package.  The important checks are these.
\begin{enumerate}[leftmargin=7mm,itemsep=3pt]
\item Every $T_h$ is odd, including $T_0=1$; every difference used as
an exponent in lifting is positive and even.
\item For odd $p\mid q$, $p\nmid q-2$.  The initial difference
$T_r-1$ therefore has valuation zero at $p$.
\item At $2$, the initial valuation is $u$, not zero; each lift adds
$u+v-1$, not merely $v$.
\item The exceptional case where $q$ is a power of $2$ is treated
separately when proving exact radix-$q$ precision.
\item The leading-difference binomial estimate uses
$2q^h\mid D_h$, rather than only $q^h\mid D_h$.
\item Exponent reduction is justified by an explicit period identity.
The code does not assume that arbitrary modular exponentiation is a
function of the exponent modulo the output modulus.
\item A repeated canonical residue $R_m=R_{m+1}$ is not confused with
$T_m\equiv T_{m+1}\pmod{q^{m+1}}$.
\item The analytic formula keeps the local sign at $2$ separate;
the higher-arrow bound excludes the exceptional base $a=2$.
\end{enumerate}
The detailed companion file \texttt{PROOF\_AUDIT.md} records the same
checks and separates mathematical proof from finite verification.

\section{Package and reproduction instructions}
The archive contains the complete \LaTeX{} source, this PDF, two
standard-library Python files, generated CSV and JSON data, a source
record, and build instructions. 

From the extracted directory, run
\begin{lstlisting}[language=bash]
python code/verify.py
python code/tower_digits.py 4 3 --height 2
python code/tower_digits.py 4 3 --height 3
python code/tower_digits.py 10 60 --rank 1000000 --argument 2
latexmk -pdf -interaction=nonstopmode -halt-on-error article.tex
\end{lstlisting}
The two finite-height commands print $27$ and $59$.  The third prints
the unpadded decimal representative for the example in
\cref{sec:arrows}.  The verification script regenerates its report,
the residue tables, the local Lambert table, and the counterexample
certificate.  Elapsed time is machine-dependent; the mathematical
outputs and assertion counts are deterministic.

\begin{thebibliography}{9}
\bibitem{oeis}
Davis Smith and OEIS contributors,
\emph{A324017: Square array of residues of $(2n-1)\uparrow\uparrow m$
modulo $(2n)^m$}, The On-Line Encyclopedia of Integer Sequences,
entry introduced 28 March 2019.
\url{https://oeis.org/A324017}; internal record:
\url{https://oeis.org/A324017/internal}.
Retrieved 19 September 2026; the internal entry showed revision
\texttt{\#91, Feb 16 2025} and retained the three conjectural comments.

\bibitem{ripa}
Marco Rip\`a,
\emph{The congruence speed formula},
arXiv:2208.02622, 2022; journal reference:
\emph{Notes on Number Theory and Discrete Mathematics} 27(4), 43--61 (2021).
\url{https://arxiv.org/abs/2208.02622}.
Used for the established decimal congruence-speed context, not as a
source of a proof of the selected OEIS statements.

\bibitem{hittmeir}
Markus Hittmeir,
\emph{A reduction of integer factorization to modular tetration},
arXiv:1707.04919, version 3, 2018;
\emph{International Journal of Foundations of Computer Science}, 2020,
DOI \texttt{10.1142/S0129054120500197}.
\url{https://arxiv.org/abs/1707.04919}.

\bibitem{mezo}
Istv\'an Mez\H{o},
\emph{The $p$-adic Lambert $W$ function},
arXiv:1801.00657, 2018.
\url{https://arxiv.org/abs/1801.00657}.
Used for the established definition and convergence theory of $W_p$.
\end{thebibliography}
\end{document}
